% GENERATED FILE. Edit canonical lessons, then run npm run generate:books.
% canonical-source-hash: fnv1a64:63f39104836baa10
% canonical-lessons: RU-C117-rabbit, RU-C117-tiger, RU-C117-butterfly, RU-C117-frog, RU-C117-tortoise

\chapter{Rabbit, Tiger, Butterfly, Frog, Tortoise}
\label{ch:ru-rabbit-tiger-butterfly-frog-tortoise}
\hlchaptermodality{117}

\begin{chapteropening}
\textbf{By the end of this chapter:} \emph{I can say a rabbit, a tiger, a butterfly, a frog and a tortoise.}

Complete the last lesson of chapter 117: I can say a rabbit, a tiger, a butterfly, a frog and a tortoise.
\end{chapteropening}

\section[\texorpdfstring{\ru{кролик} (królik)}{кролик (królik)}]{\ru{кролик} (królik) — a rabbit}

\label{lesson:RU-C117-rabbit}



\begin{quote}
Before the new one: say the Russian for a waterfall, then the Russian for a donkey.
\end{quote}

\subsection*{You'll want to know: \ru{кролик}}
\textbf{\ru{кролик}} (\emph{królik}) — \textquotedblleft{}a rabbit\textquotedblright{}.

\begin{cousinweb}
\textbf{\ru{кролик}} comes from Polish \emph{królik}, \textquotedblleft{}little king\textquotedblright{}, which copies a German dialect \emph{Künigl}, Middle High German \emph{küniklîn}, from Latin \emph{cunīculus}, \textquotedblleft{}rabbit\textquotedblright{}, heard as \textquotedblleft{}little king\textquotedblright{} (Vasmer).
\end{cousinweb}

\begin{grammarlens}[title={What you've built}]
One more word for this chapter.
\end{grammarlens}

\subsection*{Your turn}
\begin{itemize}
\raggedright
  \item \emph{Say it:} \emph{królik}
  \item \emph{Say it:} \emph{królik}, once more
  \item \emph{Say it:} say \emph{królik}
  \item \emph{Recall:} say the Russian for a waterfall, then the Russian for a donkey, then say \emph{królik} again
\end{itemize}

\begin{tcolorbox}[breakable,colback=teal!4,colframe=teal!35!black,title={Before you move on}]
What does \ru{кролик} mean? (\textquotedblleft{}A rabbit\textquotedblright{}.) Say it once more.
\end{tcolorbox}
\section[\texorpdfstring{\ru{тигр} (tigr)}{тигр (tigr)}]{\ru{тигр} (tigr) — a tiger}

\label{lesson:RU-C117-tiger}



\begin{quote}
Before the new one: say the Russian for a donkey, then the Russian for a rabbit.
\end{quote}

\subsection*{You'll want to know: \ru{тигр}}
\textbf{\ru{тигр}} (\emph{tigr}) — \textquotedblleft{}a tiger\textquotedblright{}.

\begin{cousinweb}
\textbf{\ru{тигр}} — Old Church Slavonic already had \emph{tigrŭ}, from Greek \emph{tígris}, which Vasmer traces to Avestan \emph{tighri-}, \textquotedblleft{}an arrow\textquotedblright{}; the modern word came again through French \emph{tigre} or German \emph{Tiger}, Latin \emph{tigris}.
\end{cousinweb}

\begin{grammarlens}[title={What you've built}]
One more word for this chapter.
\end{grammarlens}

\subsection*{Your turn}
\begin{itemize}
\raggedright
  \item \emph{Say it:} \emph{tigr}
  \item \emph{Say it:} \emph{tigr}, once more
  \item \emph{Say it:} say \emph{tigr}
  \item \emph{Recall:} say the Russian for a donkey, then the Russian for a rabbit, then say \emph{tigr} again
\end{itemize}

\begin{tcolorbox}[breakable,colback=teal!4,colframe=teal!35!black,title={Before you move on}]
What does \ru{тигр} mean? (\textquotedblleft{}A tiger\textquotedblright{}.) Say it once more.
\end{tcolorbox}
\section[\texorpdfstring{\ru{бабочка} (bábochka)}{бабочка (bábochka)}]{\ru{бабочка} (bábochka) — a butterfly}

\label{lesson:RU-C117-butterfly}



\begin{quote}
Before the new one: say the Russian for a rabbit, then the Russian for a tiger.
\end{quote}

\subsection*{You'll want to know: \ru{бабочка}}
\textbf{\ru{бабочка}} (\emph{bábochka}) — \textquotedblleft{}a butterfly\textquotedblright{}.

\begin{cousinweb}
\textbf{\ru{бабочка}} is a diminutive of \emph{\ru{бабка}}, \textquotedblleft{}granny\textquotedblright{} (Vasmer), from the belief that the soul of a dead person lives on as a butterfly; Greek \emph{psych\'{\={e}}} is both \textquotedblleft{}soul\textquotedblright{} and \textquotedblleft{}butterfly\textquotedblright{}.
\end{cousinweb}

\begin{grammarlens}[title={What you've built}]
One more word for this chapter.
\end{grammarlens}

\subsection*{Your turn}
\begin{itemize}
\raggedright
  \item \emph{Say it:} \emph{bábochka}
  \item \emph{Say it:} \emph{bábochka}, once more
  \item \emph{Say it:} say \emph{bábochka}
  \item \emph{Recall:} say the Russian for a rabbit, then the Russian for a tiger, then say \emph{bábochka} again
\end{itemize}

\begin{tcolorbox}[breakable,colback=teal!4,colframe=teal!35!black,title={Before you move on}]
What does \ru{бабочка} mean? (\textquotedblleft{}A butterfly\textquotedblright{}.) Say it once more.
\end{tcolorbox}
\section[\texorpdfstring{\ru{лягушка} (lyagúshka)}{лягушка (lyagúshka)}]{\ru{лягушка} (lyagúshka) — a frog}

\label{lesson:RU-C117-frog}



\begin{quote}
Before the new one: say the Russian for a tiger, then the Russian for a butterfly.
\end{quote}

\subsection*{You'll want to know: \ru{лягушка}}
\textbf{\ru{лягушка}} (\emph{lyagúshka}) — \textquotedblleft{}a frog\textquotedblright{}.

\begin{cousinweb}
\textbf{\ru{лягушка}} comes from \emph{\ru{ляга}}, \textquotedblleft{}a thigh\textquotedblright{}, and \emph{\ru{лягать}}, \textquotedblleft{}to kick\textquotedblright{} (Vasmer): the frog as the kicker.
\end{cousinweb}

\begin{grammarlens}[title={What you've built}]
One more word for this chapter.
\end{grammarlens}

\subsection*{Your turn}
\begin{itemize}
\raggedright
  \item \emph{Say it:} \emph{lyagúshka}
  \item \emph{Say it:} \emph{lyagúshka}, once more
  \item \emph{Say it:} say \emph{lyagúshka}
  \item \emph{Recall:} say the Russian for a tiger, then the Russian for a butterfly, then say \emph{lyagúshka} again
\end{itemize}

\begin{tcolorbox}[breakable,colback=teal!4,colframe=teal!35!black,title={Before you move on}]
What does \ru{лягушка} mean? (\textquotedblleft{}A frog\textquotedblright{}.) Say it once more.
\end{tcolorbox}
\section[\texorpdfstring{\ru{черепаха} (cherepákha)}{черепаха (cherepákha)}]{\ru{черепаха} (cherepákha) — a tortoise, a turtle}

\label{lesson:RU-C117-tortoise}



\begin{quote}
Before the new one: say the Russian for a butterfly, then the Russian for a frog.
\end{quote}

\subsection*{You'll want to know: \ru{черепаха}}
\textbf{\ru{черепаха}} (\emph{cherepákha}) — \textquotedblleft{}a tortoise, a turtle\textquotedblright{}.

\begin{cousinweb}
\textbf{\ru{черепаха}} is built on \emph{\ru{череп}}, \textquotedblleft{}a skull, a shell\textquotedblright{} (Vasmer), and replaced an older Slavic \emph{*žely}, \textquotedblleft{}tortoise\textquotedblright{}, for taboo reasons, after Zelenin.
\end{cousinweb}

\begin{grammarlens}[title={What you've built}]
One more word for this chapter.
\end{grammarlens}

\subsection*{Your turn}
\begin{itemize}
\raggedright
  \item \emph{Say it:} \emph{cherepákha}
  \item \emph{Say it:} \emph{cherepákha}, once more
  \item \emph{Say it:} say \emph{cherepákha}
  \item \emph{Recall:} say the Russian for a butterfly, then the Russian for a frog, then say \emph{cherepákha} again
\end{itemize}

\begin{tcolorbox}[breakable,colback=teal!4,colframe=teal!35!black,title={Before you move on}]
What does \ru{черепаха} mean? (\textquotedblleft{}A tortoise, a turtle\textquotedblright{}.) Say it once more.
\end{tcolorbox}
